Verification of LiionTR focuses on determining whether the implemented
numerical models reproduce their intended mathematical formulations.

This distinction is important:

\begin{itemize}
    \item \emph{verification} evaluates whether the equations are solved
          and implemented correctly;
    \item \emph{validation} evaluates whether those equations reproduce
          physical experimental behavior.
\end{itemize}

The current development stage places primary emphasis on verification.


\subsection{Unit-level verification}

Individual LiionTR components are tested independently wherever
possible.

The test suite includes checks for:

\begin{itemize}
    \item geometric quantities;
    \item material properties;
    \item cell mass and thermal capacity;
    \item kinetic-rate calculations;
    \item reaction-progress laws;
    \item reaction heat generation;
    \item reaction-network aggregation;
    \item gas inventory properties;
    \item gas-generation rates;
    \item ideal-gas pressure;
    \item choked and unchoked vent flow;
    \item event handling;
    \item coupled numerical simulations.
\end{itemize}

Unit-level testing is particularly important because the modular
architecture permits most physical components to be evaluated without
running a complete thermal runaway simulation.


\subsection{Dimensional consistency}

All internal physical quantities are represented in SI units.

Typical units include

\begin{align}
    T &: \si{\kelvin},
    \\
    P &: \si{\pascal},
    \\
    E_a &: \si{\joule\per\mole},
    \\
    \Delta H &: \si{\joule\per\kilogram},
    \\
    W &: \si{\kilogram\per\meter\cubed},
    \\
    \dot{Q} &: \si{\watt},
    \\
    n &: \si{\mole}.
\end{align}

Literature parameters reported in alternative units are explicitly
converted before construction of the corresponding LiionTR models.


\subsection{Reaction-energy consistency}

For a reaction with effective mass fraction $w$, reaction enthalpy
$\Delta H$, and conversion change

\begin{equation}
    \Delta \alpha
    =
    \alpha_f-\alpha_0,
\end{equation}

the expected specific energy release relative to total cell mass is

\begin{equation}
    q_{\mathrm{expected}}
    =
    w
    \Delta H
    \Delta\alpha.
    \label{eq:verification_expected_energy}
\end{equation}

For a complete reaction network,

\begin{equation}
    q_{\mathrm{expected,total}}
    =
    \sum_i
    w_i
    \Delta H_i
    \Delta\alpha_i,
    \label{eq:verification_total_energy}
\end{equation}

with the appropriate extension for multi-channel reactions.

Numerically integrating the heat-generation rate provides

\begin{equation}
    q_{\mathrm{integrated}}
    =
    \frac{1}{m_{\mathrm{cell}}}
    \int
    \dot{Q}_{\mathrm{gen}}
    \,dt.
    \label{eq:verification_integrated_energy}
\end{equation}

Agreement between
\cref{eq:verification_total_energy,eq:verification_integrated_energy}
provides an important consistency check between reaction progress and
heat generation.

\subsection{Thermal energy-balance closure}

A global energy-balance check was performed for the Hu et al.\
reference simulation.

For the lumped thermal formulation, integration over the simulation
interval gives

\begin{equation}
    q_{\mathrm{reaction}}
    =
    \Delta u_{\mathrm{sensible}}
    +
    q_{\mathrm{conv}}
    +
    \varepsilon_E,
\end{equation}

where $\varepsilon_E$ represents the numerical energy-balance residual.

The resulting balance is given in
\cref{tab:hu_energy_balance}.

\input{generated/hu2020_energy_balance}



A small residual demonstrates consistency between the integrated
reaction heat generation, the temperature evolution, and the
convective heat transfer computed by the lumped thermal model.

Because the reaction and convective energies are reconstructed by
numerical quadrature from the adaptive solver output, the residual
reflects both integration error and post-processing quadrature error.
It should therefore be interpreted as a numerical consistency metric
rather than as an independent physical validation.

The cumulative energy balance is shown in
\cref{fig:hu_cumulative_energy_balance}.

\begin{figure}[htbp]
    \centering
    \includegraphics[
        width=0.82\textwidth
    ]{hu2020_energy_balance_cumulative.pdf}

    \caption{
        Cumulative energy-balance closure for the Hu et al.\
        reference simulation. The cumulative reaction energy is
        compared with the sum of the sensible-energy increase and
        convective heat loss.
    }
    \label{fig:hu_cumulative_energy_balance}
\end{figure}

The near-superposition of the cumulative reaction energy and the
accounted thermal energy provides a direct visual verification of the
implemented lumped thermal balance.

\subsection{Hu model energy check}

A representative numerical simulation of the implemented Hu et al.\
reaction network was performed from an elevated initial temperature in
order to activate the thermal runaway mechanism.

For the tested case, the integrated specific reaction energy was

\begin{equation}
    q_{\mathrm{integrated}}
    \approx
    \SI{386.6}{\kilo\joule\per\kilogram}.
\end{equation}

The integrated heat release was consistent with the energy obtained from
the change in reaction states and the corresponding reaction
enthalpies.

This test verifies the internal consistency of the implemented reaction
network and heat-generation calculation. It should not be interpreted as
experimental validation of the Hu model.


\subsection{Thermal-response consistency}

For an adiabatic system with constant heat capacity and no heat loss,

\begin{equation}
    C_{\mathrm{th}}
    \frac{dT}{dt}
    =
    \dot{Q}_{\mathrm{gen}}.
\end{equation}

Integrating gives

\begin{equation}
    \Delta T
    =
    \frac{
        Q_{\mathrm{gen}}
    }{
        C_{\mathrm{th}}
    }.
    \label{eq:verification_adiabatic_temperature}
\end{equation}

This relationship provides a direct verification target for simplified
thermal problems.


\subsection{Gas-generation conservation}

For a closed system with no venting, the numerically integrated gas
inventory should agree with the cumulative reaction-yield relation.

For species $s$,

\begin{equation}
    n_s(t)
    -
    n_s(0)
    =
    \sum_i
    Y_{i,s}
    m_{\mathrm{cell}}
    w_i
    \left[
        \alpha_i(t)-\alpha_i(0)
    \right].
    \label{eq:verification_gas_conservation}
\end{equation}

Agreement verifies consistency between reaction progress, gas-generation
rates, and integrated species inventories.


\subsection{Pressure verification}

For a closed ideal-gas problem, pressure reconstructed from the numerical
state must satisfy

\begin{equation}
    P
    =
    \frac{
        n_{\mathrm{tot}}RT
    }{
        V_{\mathrm{free}}
    }.
\end{equation}

Synthetic problems with constant temperature or constant gas amount can
therefore be compared directly against analytical pressure evolution.


\subsection{Vent-flow verification}

The venting implementation is tested separately in both compressible-flow
regimes.

For

\begin{equation}
    \frac{P_d}{P_u}
    \leq
    r_{\mathrm{crit}},
\end{equation}

the choked-flow expression must be selected.

For

\begin{equation}
    \frac{P_d}{P_u}
    >
    r_{\mathrm{crit}},
\end{equation}

the unchoked-flow expression must be selected.

The model also verifies that

\begin{equation}
    P_u
    \leq
    P_d
\end{equation}

produces zero outward discharge.


\subsection{Vented coupled simulation}

A synthetic coupled simulation has been used to verify the complete
sequence

\begin{equation}
    \text{reaction}
    \rightarrow
    \text{heat}
    \rightarrow
    T
    \rightarrow
    \text{gas generation}
    \rightarrow
    P
    \rightarrow
    \text{vent opening}
    \rightarrow
    \text{gas discharge}.
\end{equation}

The test demonstrates that:

\begin{enumerate}
    \item reaction progress generates heat;
    \item heat generation increases temperature;
    \item reaction progress generates gas;
    \item gas accumulation and heating increase pressure;
    \item the vent opens when the prescribed pressure is reached;
    \item the solver transitions irreversibly to the open-vent system;
    \item gas discharge reduces the internal inventory;
    \item pressure approaches the downstream pressure after the dominant
          gas-generation period.
\end{enumerate}

This synthetic problem is designed to verify model coupling and event
handling rather than reproduce a particular commercial battery cell.


\subsection{Regression testing}

Automated regression tests are used to detect unintended changes in
physical behavior as the codebase evolves.

This is particularly important for coupled models because a local
software modification can alter:

\begin{itemize}
    \item reaction rates;
    \item integrated reaction energy;
    \item event times;
    \item maximum temperature;
    \item pressure evolution;
    \item vent discharge.
\end{itemize}

Regression tests therefore complement analytical unit tests.


\subsection{Verification hierarchy}

The current verification strategy follows the hierarchy illustrated in
\cref{fig:verification_hierarchy}.

\begin{figure}[htbp]
    \centering

    \begin{tikzpicture}[
        node distance=1.2cm,
        block/.style={
            rectangle,
            rounded corners,
            draw,
            align=center,
            minimum width=5.3cm,
            minimum height=0.9cm
        },
        line/.style={
            -{Latex[length=3mm]},
            thick
        }
    ]

    \node[block] (units)
        {Individual model unit tests};

    \node[block, below=of units] (analytic)
        {Analytical consistency checks};

    \node[block, below=of analytic] (coupled)
        {Coupled synthetic simulations};

    \node[block, below=of coupled] (regression)
        {Automated regression testing};

    \node[block, below=of regression] (validation)
        {Experimental validation};

    \draw[line] (units) -- (analytic);
    \draw[line] (analytic) -- (coupled);
    \draw[line] (coupled) -- (regression);
    \draw[line] (regression) -- (validation);

    \end{tikzpicture}

    \caption{
        Verification and validation hierarchy intended for LiionTR.
        Experimental validation represents a future development stage
        following verification of the numerical implementation.
    }
    \label{fig:verification_hierarchy}
\end{figure}


\subsection{Experimental validation}

Experimental validation remains an essential future step.

Relevant validation data may include:

\begin{itemize}
    \item differential scanning calorimetry measurements;
    \item accelerating rate calorimetry measurements;
    \item cell-surface temperature histories;
    \item pressure measurements;
    \item vent-opening times;
    \item total vented gas quantities;
    \item gas-composition measurements;
    \item abuse-test thermal runaway onset temperatures.
\end{itemize}

The modular structure of LiionTR is intended to allow kinetic,
thermochemical, pressure, and venting submodels to be calibrated and
validated independently before complete-cell validation is attempted.

\subsection{Coupled gas-pressure-venting verification}

The coupled gas-generation, pressure, and venting implementation was
evaluated using a deliberately simplified synthetic problem.

The purpose of this calculation is not to represent a physically
realistic lithium-ion cell, but to exercise the complete numerical
coupling

\begin{equation}
    \text{reaction}
    \rightarrow
    \text{gas generation}
    \rightarrow
    \text{pressure increase}
    \rightarrow
    \text{vent opening}
    \rightarrow
    \text{depressurization}.
\end{equation}

The model contains a single synthetic exothermic reaction producing
\ce{CO2}. The initial gas inventory consists primarily of \ce{N2} at
atmospheric pressure. The free gas volume is
\SI{1.0e-6}{\meter\cubed}, and the vent-opening pressure is prescribed as

\begin{equation}
    P_{\mathrm{vent}}
    =
    \SI{200}{\kilo\pascal}.
\end{equation}

The vent model uses an effective opening area of
\SI{1.0e-5}{\meter\squared}, a discharge coefficient of 0.8, and a
heat-capacity ratio of 1.30.

\input{generated/synthetic_vent_results}

The simulated vent opens at

\begin{equation}
    t_{\mathrm{vent}}
    =
    \SI{\SyntheticVentOpenTime}{\second},
\end{equation}

when the internal pressure reaches approximately

\begin{equation}
    P_{\max}
    =
    \SI{\SyntheticVentPeakPressure}{\kilo\pascal}.
\end{equation}

The pressure evolution is shown in
\cref{fig:synthetic_vent_pressure}.

\begin{figure}[htbp]
    \centering
    \includegraphics[
        width=0.82\textwidth
    ]{synthetic_vent_pressure.pdf}

    \caption{
        Internal pressure during the synthetic venting verification
        problem. The vent opens when the prescribed pressure threshold
        is reached, after which the internal pressure rapidly approaches
        the downstream ambient pressure.
    }
    \label{fig:synthetic_vent_pressure}
\end{figure}

At the end of the simulation, the vented system reaches

\begin{equation}
    P_f
    =
    \SI{\SyntheticVentFinalPressure}{\kilo\pascal},
\end{equation}

which is close to the prescribed downstream pressure of
\SI{101.325}{\kilo\pascal}.

For comparison, suppressing the vent while leaving the remaining model
unchanged produces a final pressure of

\begin{equation}
    P_{f,\mathrm{closed}}
    =
    \SI{\SyntheticClosedFinalPressure}{\kilo\pascal}.
\end{equation}

The very large closed-system pressure is a consequence of the
intentionally aggressive synthetic gas-generation rate and small free
volume. It should not be interpreted as a prediction for a physical
lithium-ion cell.

The internal gas inventory is shown in
\cref{fig:synthetic_vent_inventory}.

\begin{figure}[htbp]
    \centering
    \includegraphics[
        width=0.82\textwidth
    ]{synthetic_vent_gas_inventory.pdf}

    \caption{
        Internal \ce{N2}, \ce{CO2}, and total gas inventories during
        the synthetic venting verification problem.
    }
    \label{fig:synthetic_vent_inventory}
\end{figure}

Following vent opening, the internal gas inventory decreases as gas is
discharged through the vent. The corresponding reconstructed mass-flow
rate is shown in
\cref{fig:synthetic_vent_flow}.

\begin{figure}[htbp]
    \centering
    \includegraphics[
        width=0.82\textwidth
    ]{synthetic_vent_flow.pdf}

    \caption{
        Reconstructed vent mass-flow rate during the synthetic
        gas-pressure-venting verification problem.
    }
    \label{fig:synthetic_vent_flow}
\end{figure}

The maximum reconstructed vent mass-flow rate is

\begin{equation}
    \dot{m}_{\mathrm{vent,max}}
    =
    \SI{\SyntheticVentPeakMassFlow}{\milli\gram\per\second}.
\end{equation}

The irreversible vent-state implementation is additionally verified by
the automated test suite, which requires the vent-state variable to
transition from zero to one after reaching the opening pressure and
prevents any subsequent closing transition.

\subsubsection{Compressible-flow regime transition}

The synthetic verification case also exercises both compressible-flow
regimes implemented by the vent model.

For the prescribed heat-capacity ratio of $\gamma=1.30$, the critical
pressure ratio is

\begin{equation}
    r_{\mathrm{crit}}
    =
    \SyntheticCriticalPressureRatio.
\end{equation}

With a downstream pressure of
\SI{101.325}{\kilo\pascal}, the corresponding critical upstream
pressure is

\begin{equation}
    P_{\mathrm{up,crit}}
    =
    \SI{\SyntheticCriticalUpstreamPressure}{\kilo\pascal}.
\end{equation}

At vent opening, the upstream pressure is approximately
\SI{200}{\kilo\pascal}, such that

\begin{equation}
    \frac{P_{\mathrm{down}}}{P_{\mathrm{up}}}
    =
    \frac{101.325}{200}
    \approx
    0.507
    <
    r_{\mathrm{crit}}.
\end{equation}

The discharge therefore begins in the choked-flow regime.

The choked phase begins at approximately

\begin{equation}
    t_{\mathrm{choked,start}}
    =
    \SI{\SyntheticChokedFlowStart}{\milli\second}
\end{equation}

and ends at

\begin{equation}
    t_{\mathrm{choked,end}}
    =
    \SI{\SyntheticChokedFlowEnd}{\milli\second},
\end{equation}

corresponding to a duration of only

\begin{equation}
    \Delta t_{\mathrm{choked}}
    =
    \SI{\SyntheticChokedFlowDuration}{\micro\second}.
\end{equation}

The transition occurs when the interpolated upstream pressure reaches

\begin{equation}
    P_{\mathrm{up}}
    =
    \SI{\SyntheticFlowTransitionPressure}{\kilo\pascal},
\end{equation}

in agreement with the theoretical critical pressure obtained from
\cref{eq:critical_upstream_pressure}.

The very short choked-flow interval explains the sharp initial vent-flow
peak observed in \cref{fig:synthetic_vent_flow}. Once the internal
pressure falls below the critical upstream pressure, the discharge
switches to the unchoked formulation and decreases progressively as the
cell pressure approaches the downstream pressure.

\subsubsection{Gas-composition evolution}

The evolution of the internal gas composition is shown in
\cref{fig:synthetic_vent_composition}.

\begin{figure}[htbp]
    \centering
    \includegraphics[
        width=0.82\textwidth
    ]{synthetic_vent_composition.pdf}

    \caption{
        Evolution of the internal gas mole fractions during the
        synthetic venting verification problem. The initial inventory
        contains nitrogen, whereas the synthetic reaction generates
        only carbon dioxide.
    }
    \label{fig:synthetic_vent_composition}
\end{figure}

The synthetic reaction produces only \ce{CO2}, whereas the initial gas
inventory contains \ce{N2}. Consequently, the gas mixture becomes
progressively enriched in \ce{CO2}, ultimately approaching

\begin{equation}
    x_{\ce{CO2}}
    \rightarrow
    1.
\end{equation}

This compositional evolution should not be interpreted as preferential
venting of carbon dioxide. In the present mixture model, each species is
discharged according to the instantaneous gas composition. The
enrichment in \ce{CO2} results from the imposed reaction source term,
which continuously generates \ce{CO2} while no corresponding
\ce{N2} source is present.

\subsubsection{Molar conservation}

The coupled gas-generation and venting calculation was also evaluated
through a global molar balance.

Over the complete simulation interval, conservation requires

\begin{equation}
    n_0
    +
    n_{\mathrm{generated}}
    -
    n_{\mathrm{vented}}
    -
    n_f
    =
    \varepsilon_n,
\end{equation}

where $n_0$ is the initial gas inventory,
$n_{\mathrm{generated}}$ is the cumulative amount of gas generated by
the reaction model,
$n_{\mathrm{vented}}$ is the cumulative amount discharged through the
vent, $n_f$ is the final internal gas inventory, and
$\varepsilon_n$ is the numerical molar-balance residual.

\input{generated/synthetic_gas_balance}

The resulting relative residual is approximately

\begin{equation}
    \frac{\varepsilon_n}
    {n_0+n_{\mathrm{generated}}}
    \times 100
    =
    \SyntheticGasRelativeResidual\%
\end{equation}

This small residual demonstrates numerical consistency between the gas
generation model, the internal gas inventory, and the compressible vent
flow calculation.

As with the thermal energy-balance verification, the residual includes
post-processing quadrature error. In particular, the vent-flow
integration is performed only after the discrete transition to the open
vent state in order to avoid introducing a spurious trapezoidal
contribution across the vent-opening discontinuity.